\documentclass[11pt,a4paper]{amsart}

\usepackage[margin=1in]{geometry}
\usepackage{amsmath,amssymb}
\usepackage{mathtools}
\usepackage{microtype}
\usepackage[hidelinks]{hyperref}

\newtheorem{theorem}{Theorem}
\newtheorem{lemma}[theorem]{Lemma}
\newcommand{\floor}[1]{\left\lfloor #1\right\rfloor}

\title{An Explicit Fixed-Divisor Bound for A211420}
\date{July 2026}

\begin{document}

\begin{abstract}
Let
\[
a_n=\frac{(8n)!\,n!}{(4n)!(3n)!(2n)!}
\quad\text{and}\quad
L_r=\operatorname{lcm}(1,2,\ldots,8r-1).
\]
For every $n\geq0$, $r\geq1$, and $k\in\{1,2,3\}$, we prove
\[
\prod_{i=1}^{r}(kn+i)\mid L_r^r a_n.
\]
This gives an explicit uniform version of a positive-constant divisibility
assertion recorded by Peter Bala in OEIS A211420. The weaker existence of some
constant follows from Landau's multivariate factorial-ratio criterion; the
contribution here is the explicit $L_r^r$ bound and a direct prime-power proof.
The theorem has been formalized in Lean 4 with mathlib and checked by Lean's
kernel. The constant is not claimed to be sharp.
\end{abstract}

\maketitle

\section{Introduction}

Consider the factorial ratio
\[
 a_n=\frac{(8n)!\,n!}{(4n)!(3n)!(2n)!}
 \qquad(n\geq0).
\]
Peter Bala recorded in OEIS A211420 the assertion that, for fixed $r\geq1$
and $k\in\{1,2,3\}$, some positive integer $C(k,r)$, independent of $n$,
makes $C(k,r)a_n$ divisible by the consecutive product
$\prod_{i=1}^r(kn+i)$ for every $n\geq0$~\cite{OEIS}. We give the following
explicit choice.

\begin{theorem}\label{thm:main}
Let $n\geq0$, $r\geq1$, and $k\in\{1,2,3\}$. Then $a_n$ is an integer and,
with
\[
 L_r=\operatorname{lcm}(1,2,\ldots,8r-1),
 \qquad D_{k,r}(n)=\prod_{i=1}^{r}(kn+i),
\]
one has
\[
 D_{k,r}(n)\mid L_r^r a_n.
\]
In particular, the constant $L_r^r$ is positive and uniform in $n$.
\end{theorem}

The bound is not claimed to be minimal or sharp. The integrality of $a_n$ was
already known: it is the $a=4$, $b=1$ specialization of an infinite family in
Bober's classification of integral factorial ratios~\cite{Bober2009}. The
weaker existence of some $C(k,r)$ also follows from the classical multivariate
criterion of Landau~\cite{Landau1900}; a two-variable specialization gives,
for example, the much larger choice $((8r)!)^r$. The contribution here is the
explicit refinement to $L_r^r$ and a direct proof that tracks the relevant
prime powers. Rodriguez-Villegas gives broader background on integral
factorial ratios and their hypergeometric series~\cite{RV2007}.

Two published special cases are worth noting. Sun's result, under the relevant
specialization, gives $(3n+1)\mid5a_n$~\cite{Sun2012}. Yang's theorem with
$a=4$ and $b=1$ gives $(2n+1)(2n+3)\mid99a_n$ for positive $n$~\cite{Yang2015};
its two factors form an odd progression rather than the consecutive product
in Theorem~\ref{thm:main}. These comparisons are contextual only. We make no
claim of historical priority for Theorem~\ref{thm:main}.

\section{Proof of the explicit bound}

Fix $n\geq0$, $r\geq1$, and $k\in\{1,2,3\}$, and retain the notation of
Theorem~\ref{thm:main}.

For a prime $p$, we initially interpret $v_p$ as the usual valuation on
positive rational numbers. Define
\[
 \Delta(x)=\floor{8x}+\floor{x}
 -\floor{4x}-\floor{3x}-\floor{2x}.
\]
Since $8+1-4-3-2=0$, the identity
$\lfloor d(x+1)\rfloor=\lfloor dx\rfloor+d$ for integral $d$ shows that
$\Delta$ is one-periodic.

\begin{lemma}\label{lem:landau}
The function $\Delta$ takes only the values $0$ and $1$, and its support
modulo one is exactly
\[
 \left[\frac18,\frac13\right)
 \cup\left[\frac38,\frac12\right)
 \cup\left[\frac58,\frac23\right)
 \cup\left[\frac78,1\right).
\]
Moreover, for every prime $p$ and $n\geq0$,
\[
 v_p(a_n)=\sum_{j\geq1}\Delta\!\left(\frac{n}{p^j}\right),
\]
the sum is finite, and $a_n$ is an integer.
\end{lemma}

\begin{proof}
On $[0,1)$, direct evaluation between the breakpoints of the five floor
functions gives the displayed support and shows that the only values are $0$
and $1$. In particular, $\Delta(x)\geq0$ for all $x$.

Legendre's formula, applied to each factorial in $a_n$, gives the stated
valuation identity. The sum is finite: if $p^j>8n$, then every argument of a
floor occurring in the corresponding Legendre summand is smaller than $1$,
so the summand vanishes. Thus every prime valuation of the positive rational
number $a_n$ is nonnegative. Equivalently, the denominator of the factorial
ratio divides its numerator, and hence $a_n\in\N$.
\end{proof}

For an integer $q\geq1$, put
\[
 N_q=\#\{i\in\{1,\ldots,r\}:q\mid kn+i\}.
\]

\begin{lemma}\label{lem:prime-powers}
For every prime $p$,
\[
 v_p(D_{k,r}(n))=\sum_{j\geq1}N_{p^j}.
\]
If $q\geq8r$, then $N_q\leq\Delta(n/q)$. If $q<8r$, then $N_q\leq r$.
Finally,
\[
 v_p(L_r)=\#\{j\geq1:p^j<8r\}.
\]
\end{lemma}

\begin{proof}
Each factor $kn+i$ contributes one to $N_{p^j}$ for every
$1\leq j\leq v_p(kn+i)$. Summing first over factors or first over exponents
therefore gives the first identity; this is precisely why repeated prime
powers are counted with their full multiplicity.

Suppose $q\geq8r$. Since $q>r$, two distinct terms among
$kn+1,\ldots,kn+r$ cannot both be divisible by $q$, so $N_q\leq1$. If
$q\mid kn+i$, write $kn+i=mq$ and set $\varepsilon=i/(kq)$. Then
$n/q=m/k-\varepsilon$, where $0<\varepsilon\leq1/(8k)$. According to $k$
and the residue class of $m$, the fractional part of $n/q$ lies in one of
the following intervals:
\[
\begin{array}{c|c|c}
 k & m\pmod{k} & \{n/q\} \\ \hline
 1 & 0 & [7/8,1) \\
 2 & 0,1 & [15/16,1),\ [7/16,1/2) \\
 3 & 0,1,2 & [23/24,1),\ [7/24,1/3),\ [5/8,2/3)
\end{array}
\]
Every interval is contained in the support from Lemma~\ref{lem:landau}.
The left endpoints are included, notably when $q=8r$ and $i=r$, while the
right endpoints are not attained because $i>0$. Hence a divisible term forces
$\Delta(n/q)=1$, proving $N_q\leq\Delta(n/q)$.

For $q<8r$, the bound $N_q\leq r$ is immediate. Finally, the highest power of
$p$ dividing the least common multiple of $1,\ldots,8r-1$ is the largest
$p^j<8r$. Counting its possible exponents gives the last identity.
\end{proof}

\begin{proof}[Proof of Theorem~\ref{thm:main}]
Lemma~\ref{lem:landau} proves the integrality of $a_n$ before any integer
divisibility involving it is invoked. For each prime $p$, split the identity
of Lemma~\ref{lem:prime-powers} at $8r$ and apply its two bounds:
\[
\begin{aligned}
 v_p(D_{k,r}(n))
 &=\sum_{p^j<8r}N_{p^j}+\sum_{p^j\geq8r}N_{p^j} \\
 &\leq r\,v_p(L_r)+\sum_{p^j\geq8r}\Delta(n/p^j) \\
 &\leq r\,v_p(L_r)+\sum_{j\geq1}\Delta(n/p^j) \\
 &=v_p(L_r^r a_n).
\end{aligned}
\]
The second inequality uses $\Delta\geq0$ on the omitted small powers. Thus
the valuation of $D_{k,r}(n)$ does not exceed that of $L_r^r a_n$ at any
prime, and prime-by-prime comparison proves the desired divisibility. Since
$r\geq1$, the least common multiple $L_r$ is a positive integer, and therefore
$L_r^r>0$.
\end{proof}

\section{Formal verification and provenance}

The primary Lean development first proves the division-free theorem
\texttt{A211420.strengthened\_cross}. Exact factorial-ratio divisibility is
established by \texttt{A211420.integral\_cross} before the natural-number
quotient is introduced, and the quotient identity is then used to derive the
final theorem \texttt{A211420.strengthened}. The repository pins both Lean and
mathlib. Running \texttt{./scripts/verify.sh} rebuilds the development and
rejects forbidden proof escapes in project Lean files; Lean's kernel checks
the formal proof.

A separate second machine-generated Lean formalization is preserved in the
repository. It is not described as independently prompted or as human review.

The ordinary-language proof of the explicit $L_r^r$ strengthening was
generated with OpenAI GPT-5.6 Pro. OpenAI Codex produced the primary Lean 4
formalization and repository infrastructure; Lean's kernel checked the
development in the pinned environment. Danny initiated, curated, validated,
and published the project. No claim of historical priority or human peer
review is made.

\begin{thebibliography}{9}

\bibitem{Landau1900}
E. Landau,
\emph{Sur les conditions de divisibilit\'e d'un produit de factorielles par
un autre},
Nouvelles annales de math\'ematiques, 3e s\'erie, \textbf{19} (1900),
344--362.

\bibitem{Bober2009}
J. W. Bober,
\emph{Factorial ratios, hypergeometric series, and a family of step
functions},
J. London Math. Soc. (2) \textbf{79} (2009), 422--444.

\bibitem{RV2007}
F. Rodriguez-Villegas,
\emph{Integral ratios of factorials and algebraic hypergeometric functions},
arXiv:math/0701362 (2007).

\bibitem{Sun2012}
Z.-W. Sun,
\emph{On divisibility of binomial coefficients},
J. Aust. Math. Soc. \textbf{93} (2012), 189--201.

\bibitem{Yang2015}
Q.-H. Yang,
\emph{Proof of a conjecture of Amdeberhan and Moll on a divisibility
property of binomial coefficients},
Electron. J. Combin. \textbf{22} (2015), Paper P1.9.

\bibitem{OEIS}
OEIS Foundation Inc.,
\emph{The On-Line Encyclopedia of Integer Sequences},
Sequence A211420, \url{https://oeis.org/A211420}.

\end{thebibliography}

\end{document}
